% Options for packages loaded elsewhere
\PassOptionsToPackage{unicode}{hyperref}
\PassOptionsToPackage{hyphens}{url}
%
\documentclass[
]{article}
\usepackage{lmodern}
\usepackage{amssymb,amsmath}
\usepackage{ifxetex,ifluatex}
\ifnum 0\ifxetex 1\fi\ifluatex 1\fi=0 % if pdftex
  \usepackage[T1]{fontenc}
  \usepackage[utf8]{inputenc}
  \usepackage{textcomp} % provide euro and other symbols
\else % if luatex or xetex
  \usepackage{unicode-math}
  \defaultfontfeatures{Scale=MatchLowercase}
  \defaultfontfeatures[\rmfamily]{Ligatures=TeX,Scale=1}
\fi
% Use upquote if available, for straight quotes in verbatim environments
\IfFileExists{upquote.sty}{\usepackage{upquote}}{}
\IfFileExists{microtype.sty}{% use microtype if available
  \usepackage[]{microtype}
  \UseMicrotypeSet[protrusion]{basicmath} % disable protrusion for tt fonts
}{}
\makeatletter
\@ifundefined{KOMAClassName}{% if non-KOMA class
  \IfFileExists{parskip.sty}{%
    \usepackage{parskip}
  }{% else
    \setlength{\parindent}{0pt}
    \setlength{\parskip}{6pt plus 2pt minus 1pt}}
}{% if KOMA class
  \KOMAoptions{parskip=half}}
\makeatother
\usepackage{xcolor}
\IfFileExists{xurl.sty}{\usepackage{xurl}}{} % add URL line breaks if available
\IfFileExists{bookmark.sty}{\usepackage{bookmark}}{\usepackage{hyperref}}
\hypersetup{
  hidelinks,
  pdfcreator={LaTeX via pandoc}}
\urlstyle{same} % disable monospaced font for URLs
\usepackage{longtable,booktabs}
% Correct order of tables after \paragraph or \subparagraph
\usepackage{etoolbox}
\makeatletter
\patchcmd\longtable{\par}{\if@noskipsec\mbox{}\fi\par}{}{}
\makeatother
% Allow footnotes in longtable head/foot
\IfFileExists{footnotehyper.sty}{\usepackage{footnotehyper}}{\usepackage{footnote}}
\makesavenoteenv{longtable}
\setlength{\emergencystretch}{3em} % prevent overfull lines
\providecommand{\tightlist}{%
  \setlength{\itemsep}{0pt}\setlength{\parskip}{0pt}}
\setcounter{secnumdepth}{-\maxdimen} % remove section numbering

\author{}
\date{}

\begin{document}

\hypertarget{constraint-based-selection-of-gauge-structure-from-ell1-topology}{%
\section{\texorpdfstring{Constraint-Based Selection of Gauge Structure
from \(\ell^1\)
Topology}{Constraint-Based Selection of Gauge Structure from \textbackslash ell\^{}1 Topology}}\label{constraint-based-selection-of-gauge-structure-from-ell1-topology}}

\hypertarget{a-minimal-consistency-framework-for-the-emergence-of-su3-su2-u1}{%
\subsection{A Minimal Consistency Framework for the Emergence of SU(3) ×
SU(2) ×
U(1)}\label{a-minimal-consistency-framework-for-the-emergence-of-su3-su2-u1}}

\textbf{Author:} Jeremy H. Carroll \textbf{Version:} v1.0 (Pre-Print)
\textbf{Date:} March 2026

\begin{center}\rule{0.5\linewidth}{0.5pt}\end{center}

\hypertarget{abstract}{%
\subsection{Abstract}\label{abstract}}

We present a constraint-based derivation of gauge structure arising from
minimal \(\ell^1\) defect geometry on discrete systems.

Starting from an additive \(\ell^1\) inconsistency functional, we impose
a single structural principle: defect resolution must proceed via
monotone reduction of local inconsistency. This induces a minimal closed
consistency loop of size \(N = 3\), equipped with a cyclic transport
operator satisfying \(S^3 = I\).

Diagonalization of this operator forces complex phase structure via
roots of unity, while invariance under both cyclic transport and
independent phase evolution uniquely selects the \(\ell^2\) norm as the
compatible transport metric. As a consequence, admissible dynamics are
unitary without being assumed.

We then classify the connected compact symmetry groups acting
irreducibly on the resulting state spaces under these constraints. For
the minimal closed loop, this yields \(SU(3)\) as the unique admissible
non-abelian symmetry under the stated structural constraints.
Lower-dimensional projections of the same defect-resolution mechanism
generate \(SU(2)\) from binary ambiguity in non-unique extensions and
\(U(1)\) from one-dimensional phase transport.

The resulting structure reproduces the Standard Model gauge group:
\[SU(3) \times SU(2) \times U(1)\] as a consequence of minimal
consistency requirements, rather than symmetry breaking from a higher
unified group.

This framework provides a structural selection mechanism for gauge
symmetry independent of parameter fitting. While the physical
realization of these constraints remains an open question, the
mathematical result establishes a unique compatibility class between
\(\ell^1\) defect geometry, cyclic transport, and unitary symmetry.

\begin{center}\rule{0.5\linewidth}{0.5pt}\end{center}

\hypertarget{introduction}{%
\subsection{1. Introduction}\label{introduction}}

Grand Unified Theories have traditionally assumed that the Standard
Model gauge group \(SU(3) \times SU(2) \times U(1)\) is the low-energy
remnant of a larger simple group---\(SU(5)\) {[}1{]}, \(SO(10)\)
{[}2{]}, or an exceptional group---broken at high energies. This
assumption leads to predictions (proton decay, magnetic monopoles, large
unified representations) that have not been experimentally confirmed,
and in some cases have been excluded {[}3{]}.

This work takes a different approach. The Standard Model gauge group may
not be a broken symmetry of a larger theory, but the unique structure
compatible with minimal consistency on a discrete space. We do not claim
to deduce the necessity of the Standard Model from pure logic; instead,
given explicit geometric constraints, we show the Standard Model gauge
group is uniquely selected directly from the \(\ell^1\) topological
structure established in Paper 008. The foundational requirement is the
same: local corrections to inconsistency must reduce, not increase,
global inconsistency. From this single axiom, the 2-simplex is forced,
and its dimensional boundaries uniquely determine the gauge structure.

\textbf{Scope.} This paper derives the gauge group, generation count,
Weinberg angle, and electroweak breaking pattern. It does not derive the
fermion mass spectrum, Yukawa couplings, or CKM/PMNS mixing matrices.
These remain open problems.

\begin{center}\rule{0.5\linewidth}{0.5pt}\end{center}

\hypertarget{core-derivation-from-ell1-obstruction-to-gauge-structure}{%
\subsection{\texorpdfstring{2. Core Derivation: From \(\ell^1\)
Obstruction to Gauge
Structure}{2. Core Derivation: From \textbackslash ell\^{}1 Obstruction to Gauge Structure}}\label{core-derivation-from-ell1-obstruction-to-gauge-structure}}

\hypertarget{ell1-defect-geometry-and-minimal-closure}{%
\subsubsection{\texorpdfstring{2.1 \(\ell^1\) Defect Geometry and
Minimal
Closure}{2.1 \textbackslash ell\^{}1 Defect Geometry and Minimal Closure}}\label{ell1-defect-geometry-and-minimal-closure}}

We begin with a defect functional defined on discrete configurations:
\[|\delta|_1 = \sum_e |\delta_e|\] representing the total inconsistency
across local relations.

A fundamental constraint is \textbf{monotone defect resolution}: any
admissible refinement must not increase \(\ell^1\) inconsistency.

Starting from a minimal contradiction on two nodes, iterative
subdivision produces a configuration with three mutually interacting
components. This yields: \[N = 3\] as the minimal nontrivial closed
consistency loop.

\hypertarget{cyclic-transport-and-emergence-of-complex-structure}{%
\subsubsection{2.2 Cyclic Transport and Emergence of Complex
Structure}\label{cyclic-transport-and-emergence-of-complex-structure}}

The minimal closed structure admits a cyclic shift operator:
\[S : \mathbb{C}^3 \to \mathbb{C}^3, \quad S^3 = I\]

Diagonalizing \(S\) yields eigenvalues:
\[\lambda = 1, e^{2\pi i/3}, e^{-2\pi i/3}\]

Thus: * cyclic transport forces roots of unity * real representations
fail to decouple modes * complex scalars are required

\(\implies\) state space is intrinsically \(\mathbb{C}^3\).

\hypertarget{emergence-of-the-invariant-transport-metric}{%
\subsubsection{2.3 Emergence of the Invariant Transport
Metric}\label{emergence-of-the-invariant-transport-metric}}

The \(\ell^1\) functional defines defect magnitude but is not preserved
under diagonalization of \(S\).

We seek a norm satisfying: * invariance under cyclic transport *
invariance under independent phase evolution in the diagonal basis

These constraints restrict admissible norms to those depending only on
mode magnitudes, symmetric under permutation, and additive across
independent modes. Under these conditions, the \(\ell^2\) norm is the
unique norm compatible with linear decomposition into independent
spectral modes. A formal proof of this uniqueness is provided in the
companion derivation.

\hypertarget{unitary-evolution}{%
\subsubsection{2.4 Unitary Evolution}\label{unitary-evolution}}

Let \(T\) be a linear evolution operator. Metric preservation implies:
\[|Tx|_2 = |x|_2\]

Thus: \[T^\dagger T = I \implies G \subset U(3)\]

Unitary structure is therefore not assumed, but forced by symmetry.

\hypertarget{su3-from-irreducible-cyclic-closure}{%
\subsubsection{\texorpdfstring{2.5 \(SU(3)\) from Irreducible Cyclic
Closure}{2.5 SU(3) from Irreducible Cyclic Closure}}\label{su3-from-irreducible-cyclic-closure}}

We now classify admissible symmetry groups acting on \(\mathbb{C}^3\).

Impose: * irreducibility (no invariant subspaces) * compatibility with
cyclic transport (commutation with \(S\)) * nontrivial holonomy
(non-abelian structure) * invariance under global phase

Then: * irreducibility excludes product groups * global phase acts
trivially \(\implies\) projective symmetry * determinant constraint
yields: \[G = SU(3)\]

This identifies \(SU(3)\) as the unique connected compact Lie group
satisfying these constraints within the class of irreducible unitary
representations on \(\mathbb{C}^3\). A full classification argument is
provided in Appendix A.

\hypertarget{su2-from-minimal-ambiguity}{%
\subsubsection{\texorpdfstring{2.6 \(SU(2)\) from Minimal
Ambiguity}{2.6 SU(2) from Minimal Ambiguity}}\label{su2-from-minimal-ambiguity}}

Not all defects admit unique \(\ell^1\)-consistent resolution. In the
minimal nontrivial case, exactly two inequivalent extensions exist.

Thus: state space \(\cong \mathbb{C}^2\).

Repeating the previous structure: * symmetry-preserving transformations
form \(U(2)\) * global phase is unobservable

\[G = U(2)/U(1) \cong SU(2)\]

Thus \(SU(2)\) arises from \textbf{binary ambiguity}, not assumption.

\hypertarget{u1-from-edge-transport}{%
\subsubsection{\texorpdfstring{2.7 \(U(1)\) from Edge
Transport}{2.7 U(1) from Edge Transport}}\label{u1-from-edge-transport}}

Restricting to a single directed edge: * state space is one-dimensional
* transformations preserve magnitude

Thus: \[T = e^{i\theta} \implies G = U(1)\]

\hypertarget{unified-gauge-structure}{%
\subsubsection{2.8 Unified Gauge
Structure}\label{unified-gauge-structure}}

The three structures arise as projections of the same defect-resolution
mechanism:

\begin{longtable}[]{@{}lll@{}}
\toprule
Structure & Origin & Symmetry\tabularnewline
\midrule
\endhead
3-cycle closure & minimal loop & \(SU(3)\)\tabularnewline
2-way ambiguity & non-unique extension & \(SU(2)\)\tabularnewline
edge transport & phase evolution & \(U(1)\)\tabularnewline
\bottomrule
\end{longtable}

Thus: \[G = SU(3) \times SU(2) \times U(1)\]

\hypertarget{epistemic-boundary}{%
\subsubsection{2.9 Epistemic Boundary}\label{epistemic-boundary}}

This derivation establishes: * a structural mechanism selecting the
gauge group * independence from parameter fitting * robustness under
variation of metric details

It does \textbf{not} establish: * necessity of the initial \(\ell^1\)
constraint * derivation of mass spectra or mixing parameters *
uniqueness of physical realization

\begin{center}\rule{0.5\linewidth}{0.5pt}\end{center}

\hypertarget{exclusion-of-traditional-gut-groups}{%
\subsection{3. Exclusion of Traditional GUT
Groups}\label{exclusion-of-traditional-gut-groups}}

\hypertarget{su5-the-georgiglashow-model}{%
\subsubsection{\texorpdfstring{3.1 \(SU(5)\): The Georgi--Glashow
Model}{3.1 SU(5): The Georgi--Glashow Model}}\label{su5-the-georgiglashow-model}}

\(SU(5)\) requires a 5-dimensional fundamental representation. The
2-simplex has 3 edges, giving \(C^1 = \mathbb{C}^3\), not
\(\mathbb{C}^5\). To accommodate \(SU(5)\), one would need a 5-gon
(pentagon), but the transport amplitude
\(a(5) = \mathrm{sinc}(\pi \cdot 3/(2 \cdot 5)) = 0.858\) is strictly
less than \(a(3) = 3/\pi = 0.955\). The pentagon is suboptimal: the
\(\ell^1\) healing selects the triangle, not the pentagon.

Additionally, \(SU(5)\) predicts proton decay at rates
\(\tau_p \sim 10^{29}\)--\(10^{31}\) years, which is excluded by the
Super-Kamiokande bound \(\tau_p > 2.4 \times 10^{34}\) years {[}3{]} by
three or more orders of magnitude.

\hypertarget{so10-and-exceptional-groups}{%
\subsubsection{\texorpdfstring{3.2 \(SO(10)\) and Exceptional
Groups}{3.2 SO(10) and Exceptional Groups}}\label{so10-and-exceptional-groups}}

\(SO(10)\) requires a 10-dimensional representation
(\(a(10) = 0.757 \ll a(3)\)). \(E_6\) has the correct \(\mathbb{Z}_3\)
center but requires a 27-dimensional representation (\(a(27) = 0.683\)).
\(E_8\) has trivial center and 248-dimensional minimum faithful
representation. All fail constraint C1 (dimension \(\neq 3\)).

Within the \(\ell^1\) transport framework, higher-order polygons
(\(N > 3\)) produce strictly smaller transport amplitudes under the
chosen averaging procedure. This suggests that the triangular
configuration is optimal under the monotone defect reduction criterion.
While this does not constitute a formal impossibility proof for larger
symmetry groups, it provides a structural selection mechanism favoring
\(N = 3\) over \(N > 3\) within this framework.

\hypertarget{why-the-standard-model-is-not-broken-from-a-larger-group}{%
\subsubsection{3.3 Why the Standard Model Is Not Broken from a Larger
Group}\label{why-the-standard-model-is-not-broken-from-a-larger-group}}

Traditional GUTs assume a top-down pattern:
\(G_{\mathrm{big}} \xrightarrow{\mathrm{breaking}} G_{\mathrm{SM}}\).
The \(\ell^1\) framework shows the opposite: \(G_{\mathrm{SM}}\) is the
direct output of the minimal topology. Within the \(\ell^1\)-constrained
transport framework, no larger group satisfies the minimality and
optimality conditions. The three gauge factors arise from dimensional
projection of a single 2-simplex, not from the fragmentation of a higher
symmetry.

\begin{center}\rule{0.5\linewidth}{0.5pt}\end{center}

\hypertarget{the-coupling-constant-and-the-axiom-of-geometric-filtering}{%
\subsection{4. The Coupling Constant and the Axiom of Geometric
Filtering}\label{the-coupling-constant-and-the-axiom-of-geometric-filtering}}

\hypertarget{the-edge-coupling-and-structural-fragility}{%
\subsubsection{4.1 The Edge Coupling and Structural
Fragility}\label{the-edge-coupling-and-structural-fragility}}

From Paper 008: the geometric transport amplitude on the optimal
\(\ell^1\) simplex was evaluated as a first-order classical
approximation \(a = 3/\pi\), yielding a unitarity deficit
\(D = 1 - (3/\pi)^2 = (\pi^2 - 9)/\pi^2\), which gives an edge coupling
\(\alpha_{\mathrm{edge}}^{-1} = \pi^2/(\pi^2 - 9) \approx 11.35\).

\textbf{Epistemic Caveat:} The exact value \(a = 3/\pi\) is the most
fragile analytic step in the entire \(\ell^1\) framework. It depends
directly on the chosen angular measure, the domain integration
boundaries along the 1-dimensional phase coordinate, and a classical
scaling assumption prior to field quantization. If this specific
calculation falls to higher-order corrections or requires a superior
measure-theoretic definition, the subsequent numerical couplings will
shift.

However, the \textbf{gauge structure
(\(SU(3) \times SU(2) \times U(1)\)) and the generation count
(\(N_g=3\)) rigorously survive}. They are rigidly protected topological
invariants of the \(\ell^1\) constraints, unconditionally decoupled from
the precise numerical value of the transport amplitude. The strength of
the \(\ell^1\) framework lies in its capacity as a structural constraint
system that reproduces key aspects of the Standard Model without
parameter fitting, rather than claiming absolute numerical finality over
dynamic properties. Importantly, all group-theoretic results presented
in Sections 2--3 are independent of this numerical value.

\hypertarget{the-casimir-scaling}{%
\subsubsection{4.2 The Casimir Scaling}\label{the-casimir-scaling}}

Edge transport acts in the fundamental representation of \(SU(3)\) (dim
3); loop curvature acts in the adjoint representation (dim 8). The
quadratic Casimir invariants are \(C_2(\mathrm{fund}) = 4/3\) and
\(C_2(\mathrm{adj}) = 3\), giving the ratio
\(C_2(\mathrm{adj})/C_2(\mathrm{fund}) = 9/4\).

The GUT coupling:

\[\alpha_{\mathrm{GUT}}^{-1} = \frac{C_2(\mathrm{adj})}{C_2(\mathrm{fund})} \times \alpha_{\mathrm{edge}}^{-1} = \frac{9}{4} \times \frac{\pi^2}{\pi^2 - 9} \approx 25.54\]

This result contains zero free parameters.

\hypertarget{gauge-invariance}{%
\subsubsection{4.3 Gauge Invariance}\label{gauge-invariance}}

The trace of the unitarity deficit \(\mathrm{Tr}(D)\) is invariant under
arbitrary local unitary gauge transformations. For the cyclic loop
\(T_\Delta = T_{12} T_{23} T_{31}\), a gauge transformation
\(V(i) \in U(N)\) at vertex \(i\) gives
\(T'_\Delta = V(1) T_\Delta V(1)^\dagger\), and
\(\mathrm{Tr}(D') = \mathrm{Tr}(D)\) by the cyclic property of the
trace. Computational verification confirms
\(|\mathrm{Tr}(D') - \mathrm{Tr}(D)| < 10^{-10}\) across 1000 random
transformations (\texttt{execute\_33} in Paper 008).

\begin{center}\rule{0.5\linewidth}{0.5pt}\end{center}

\hypertarget{the-generation-count-spectral-inevitability}{%
\subsection{5. The Generation Count: Spectral
Inevitability}\label{the-generation-count-spectral-inevitability}}

\hypertarget{the-transport-operator-and-spectral-sectors}{%
\subsubsection{5.1 The Transport Operator and Spectral
Sectors}\label{the-transport-operator-and-spectral-sectors}}

The generation count \(N_g\) is not an imported parameter; it is
algebraically forced by the spectral properties of the \(\ell^1\)
transport operator on the global lattice.

The \(\ell^1\) transport operator defines a discrete convolution algebra
on the 2-simplex. Because the minimal transport solution uniquely forces
a \(\mathbb{Z}_3\) cyclic symmetry (\(N=3\)), its Fourier dual
decomposes into exactly three irreducible characters:
\(\chi_k(j) = \exp(i 2\pi k j / 3)\) for \(k \in \{0, 1, 2\}\).

In momentum space (the Brillouin zone of the \(\{3,6\}\) tessellation),
these three orthogonal representations correspond to exactly three
distinct, symmetry-protected invariant zero-modes (Dirac valleys).
Standard Nielsen-Ninomiya doubling forces chiral fermions to appear in
pairs (\(N_g \in \{2, 4, ...\}\)) on a hypercubic lattice due to spatial
inversion symmetries. We propose that the \(\ell^1\)-stabilized
triangular lattice, which explicitly breaks parity (\(\mathcal{P}\)) and
time-reversal (\(\mathcal{T}\)) at the unit cell structure (via directed
cyclic boundaries), structurally bypasses classical hypercubic pairing.
In this context, the algebraic index of the chiral transport operator is
instead governed by the \(\mathbb{Z}_3\) vertex coloring of the
\(\{3,6\}\) lattice. This construction yields three distinct symmetry
sectors, which we interpret as corresponding to the observed generation
count. A full index-theoretic derivation remains an open problem.

\hypertarget{monotonicity-and-the-exclusion-of-n-neq-3}{%
\subsubsection{\texorpdfstring{5.2 Monotonicity and the Exclusion of
\(N \neq 3\)}{5.2 Monotonicity and the Exclusion of N \textbackslash neq 3}}\label{monotonicity-and-the-exclusion-of-n-neq-3}}

We can explicitly exclude other generation counts using the same strict
monotonicity logic that excluded higher-rank gauge groups like \(SU(5)\)
and \(SO(10)\):

\begin{itemize}
\tightlist
\item
  \textbf{\(N = 1\):} A single generation cannot resolve the
  \(\ell^1\)-induced chiral defect because the defect is intrinsically a
  global cyclic obstruction. Attempting to restrict the chiral modes to
  a single uniform phase state collapses the discrete transport,
  yielding \(a(1) = 0\). The minimal consistent transport solution must
  span all three symmetry-equivalent local sectors simultaneously.
\item
  \textbf{\(N = 2\):} Two sectors lack the topological complexity to
  support a non-trivial cyclic geometric defect (a directed loop).
  \(N=2\) can only support an alternating bipartite graph structure,
  which produces exactly zero bulk flux/curvature.
\item
  \textbf{\(N > 3\):} Any generation count \(N_g > 3\) requires tiling a
  higher-order manifold as the fundamental irreducible unit. As proven
  in Section 3, the transport amplitude \(a(N)\) is strictly decreasing.
  Adding a fourth generation requires a fundamentally sub-optimal
  geometry, explicitly breaking the \(\ell^1\) action minimum.
\end{itemize}

\hypertarget{global-anomaly-cancellation}{%
\subsubsection{5.3 Global Anomaly
Cancellation}\label{global-anomaly-cancellation}}

While the standard cubic gauge anomaly \(\sum_f Q_f^3 = 0\) cancels
identically within a single generation:

\[3 \times \left(\frac{2}{3}\right)^3 + 3 \times \left(-\frac{1}{3}\right)^3 + 0^3 + (-1)^3 + 3 \times \left(-\frac{2}{3}\right)^3 + 3 \times \left(\frac{1}{3}\right)^3 + 1^3 = 0\]

a solitary generation fails the \emph{topological} anomaly test. By the
discrete Atiyah-Singer index theorem, the purely geometric index of the
Dirac operator over the triangular space must be matched by the
analytical index (the number of zero-modes). Because the geometric
stabilization requires 3 independently colored vertices to tile the
\(\ell^1\) cross-polytope without topological frustration, the
analytical index requires exactly \(N_g = 3\) distinct chiral sectors. A
single generation leaves the global lattice anomaly fatally uncancelled.

\begin{center}\rule{0.5\linewidth}{0.5pt}\end{center}

\hypertarget{the-weinberg-angle-and-electroweak-structure}{%
\subsection{6. The Weinberg Angle and Electroweak
Structure}\label{the-weinberg-angle-and-electroweak-structure}}

\hypertarget{sin2theta_w-38-at-m_mathrmgut}{%
\subsubsection{\texorpdfstring{6.1 \(\sin^2\theta_W = 3/8\) at
\(M_{\mathrm{GUT}}\)}{6.1 \textbackslash sin\^{}2\textbackslash theta\_W = 3/8 at M\_\{\textbackslash mathrm\{GUT\}\}}}\label{sin2theta_w-38-at-m_mathrmgut}}

At exact unification (\(\alpha_1 = \alpha_2 = \alpha_3\) at
\(M_{\mathrm{GUT}}\)), the Weinberg angle is determined by the
normalization of the \(U(1)_Y\) generator. Traditional GUTs import
\(k=5/3\) by embedding the Standard Model into \(SU(5)\). However, since
the \(\ell^1\) minimum strictly excludes \(SU(5)\) (Section 3), we
cannot use its embedding to fix the trace.

Instead, the normalization factor \(k=5/3\) is derived intrinsically by
evaluating the trace orthogonality of the generators across the
projected geometries. On the full \(\mathbb{C}^3\) bulk, the trace
squared of the hypercharge assignments yields
\(\mathrm{Tr}(Y^2) = 2 \times (1/2)^2 + 2 \times (1/2)^2 + 1^2 + 1^2 = 10/3\).
On the \(\mathbb{C}^2\) open hinge, the trace squared of the isospin
generator yields
\(\mathrm{Tr}(T_3^2) = 4 \times (1/2)^2 + 2 \times (1/2)^2 = 2\). The
intrinsic geometric scale factor between the \(U(1)_Y\) projection and
the \(SU(2)\) projection is exactly the ratio of these traces:
\(k = \mathrm{Tr}(Y^2)/\mathrm{Tr}(T_3^2) = (10/3) / 2 = 5/3\).

This internal representation-theoretic scaling gives:

\[\sin^2\theta_W(M_{\mathrm{GUT}}) = \frac{1}{1 + 5/3} = \frac{3}{8} = 0.375\]

This is an exact result at \(M_{\mathrm{GUT}}\). The experimental value
at \(M_Z\) (\(\sin^2\theta_W = 0.2312\)) differs due to RG running
(\texttt{execute\_44}).

\hypertarget{the-electroweak-breaking-pattern}{%
\subsubsection{6.2 The Electroweak Breaking
Pattern}\label{the-electroweak-breaking-pattern}}

The 2-simplex has a natural dimensional decomposition:

\begin{itemize}
\tightlist
\item
  \textbf{3D bulk}: \(SU(3)\) (color confinement)
\item
  \textbf{2D face}: \(SU(2) \times U(1)_Y\) (electroweak)
\item
  \textbf{1D edge}: \(U(1)_{\mathrm{em}}\) (electromagnetism)
\end{itemize}

The electroweak breaking \(SU(2) \times U(1)_Y \to U(1)_{\mathrm{em}}\)
corresponds to restricting the 2-dimensional face symmetry to the
1-dimensional edge symmetry. This is a geometric dimensional projection,
not spontaneous symmetry breaking.

The Higgs mechanism provides the mass scale (\(v = 246\) GeV) and the
mass ratios of \(W\), \(Z\), and fermions. It does not determine the
breaking pattern, which is geometric.

\hypertarget{the-charge-formula}{%
\subsubsection{6.3 The Charge Formula}\label{the-charge-formula}}

The Gell-Mann--Nishijima formula \(Q_{\mathrm{em}} = T_3 + Y/2\) is the
projection from the 2-dimensional face quantum numbers (\(T_3\) from
\(SU(2)\), \(Y\) from \(U(1)_Y\)) to the 1-dimensional edge charge
(\(Q_{\mathrm{em}}\)). All Standard Model fermion charges are reproduced
exactly (\texttt{execute\_45}).

\begin{center}\rule{0.5\linewidth}{0.5pt}\end{center}

\hypertarget{renormalization-group-flow}{%
\subsection{7. Renormalization Group
Flow}\label{renormalization-group-flow}}

\hypertarget{imported-quantities}{%
\subsubsection{7.1 Imported Quantities}\label{imported-quantities}}

The RG flow from \(\alpha_{\mathrm{GUT}}^{-1} \approx 25.54\) to
infrared observables requires:

\begin{itemize}
\tightlist
\item
  \(M_{\mathrm{GUT}} \approx 1.7 \times 10^{16}\) GeV (phenomenological
  anchor)
\item
  \(M_{\mathrm{SUSY}} \approx 800\) GeV (phenomenological anchor)
\item
  MSSM beta coefficients: \(b_1 = 33/5\), \(b_2 = 1\), \(b_3 = -3\)
  (above \(M_{\mathrm{SUSY}}\))
\item
  SM beta coefficients: \(b_1 = 41/10\), \(b_2 = -19/6\), \(b_3 = -7\)
  (below \(M_{\mathrm{SUSY}}\))
\end{itemize}

These quantities encode the matter content and are not derived from
\(\ell^1\) topology in this work.

\hypertarget{results-at-m_z}{%
\subsubsection{\texorpdfstring{7.2 Results at
\(M_Z\)}{7.2 Results at M\_Z}}\label{results-at-m_z}}

\begin{longtable}[]{@{}llll@{}}
\toprule
Observable & Prediction & PDG 2024 {[}4{]} & Deviation\tabularnewline
\midrule
\endhead
\(\alpha_{\mathrm{em}}^{-1}(M_Z)\) & \(127.97\) & \(127.951\) &
\(+0.02\%\)\tabularnewline
\(\alpha_s(M_Z)\) & \(0.1181\) & \(0.1180\) & \(+0.11\%\)\tabularnewline
\(\sin^2\theta_W(M_Z)\) & \(0.2292\) & \(0.2312\) &
\(-0.88\%\)\tabularnewline
\bottomrule
\end{longtable}

\hypertarget{falsifiable-prediction}{%
\subsubsection{7.3 Falsifiable
Prediction}\label{falsifiable-prediction}}

The two-parameter fit (\(M_{\mathrm{GUT}}\), \(M_{\mathrm{SUSY}}\))
yields an optimal \(M_{\mathrm{SUSY}} \approx 800\) GeV. Current
ATLAS/CMS exclusion bounds constrain specific MSSM spectra but do not
universally exclude this range. If no supersymmetric signatures are
observed at the HL-LHC or a future collider, the MSSM realization of
this framework is falsified.

\begin{center}\rule{0.5\linewidth}{0.5pt}\end{center}

\hypertarget{epistemic-limits-derived-vs.-imported}{%
\subsection{8. Epistemic Limits: Derived
vs.~Imported}\label{epistemic-limits-derived-vs.-imported}}

To maintain rigorous epistemic boundaries, the framework strictly
segregates what is robustly derived from what requires phenomenological
or analytical import.

\hypertarget{constraint-derived-topologically-rigid}{%
\subsubsection{8.1 Constraint-Derived (Topologically
Rigid)}\label{constraint-derived-topologically-rigid}}

\begin{enumerate}
\def\labelenumi{\arabic{enumi}.}
\tightlist
\item
  \textbf{The \(SU(3) \times SU(2) \times U(1)\) gauge group structure}
  (conditional strictly on C1--C5 constraints).
\item
  \textbf{The maximum generation count} (\(N_g = 3\)).
\item
  \textbf{The trace normalization of the \(U(1)\) and \(SU(2)\)
  projections}, yielding exactly \(\sin^2\theta_W = 3/8\) at
  unification.
\item
  \textbf{The electroweak breaking pattern} as a 2D-to-1D boundary
  projection, underlying the charge quantization formula
  (\(Q = T_3 + Y/2\)).
\item
  \textbf{Exclusion of higher-rank GUTs} (\(SU(5), SO(10), E_6, E_8\))
  via transport amplitude monotonicity.
\end{enumerate}

\hypertarget{measure-dependent-analytically-fragile}{%
\subsubsection{8.2 Measure-Dependent (Analytically
Fragile)}\label{measure-dependent-analytically-fragile}}

\begin{enumerate}
\def\labelenumi{\arabic{enumi}.}
\tightlist
\item
  \textbf{The numerical geometric transport amplitude}
  (\(a = \mathrm{sinc}(\pi/6) = 3/\pi\)). Deriving this constant relies
  on averaging \(\cos(\theta)\) over a chosen domain with a Haar
  measure. While this leads to an extremely accurate
  \(\alpha_{\mathrm{GUT}}^{-1} \approx 25.54\), the calculation
  currently lacks a uniqueness theorem for the measure choice.
\end{enumerate}

\hypertarget{ir-fitted-phenomenologically-imported}{%
\subsubsection{8.3 IR-Fitted (Phenomenologically
Imported)}\label{ir-fitted-phenomenologically-imported}}

\begin{enumerate}
\def\labelenumi{\arabic{enumi}.}
\tightlist
\item
  \textbf{The ultraviolet cut-off scale}: \(M_{\mathrm{GUT}}\).
\item
  \textbf{Explicit gauge running scales} (e.g.~\(M_{\mathrm{SUSY}}\))
  down to \(M_Z\), which require standard beta-function integration
  using phenomenological matter content.
\item
  \textbf{The Higgs mechanism}: though not responsible for the breaking
  \emph{pattern}, it remains required as the mass scale generator
  (\(v = 246 \text{ GeV}\)).
\item
  \textbf{Fermion mass spectrum and CKM/PMNS matrices}: Not derived;
  allocated as open problems.
\end{enumerate}

\hypertarget{summary-verdict}{%
\subsubsection{8.4 Summary Verdict}\label{summary-verdict}}

The Standard Model structure is demonstrated to be the unique solution
under these explicit \(\ell^1\) topological constraints. While this
mechanism beautifully selects precisely the observed physical
symmetries---explaining \emph{how} the geometry enforces empirical
phenomenology without parameter fitting---the necessity of the
constraints themselves (\emph{why} nature insists on this precise
geometric minimum) remains an open physical question.

\begin{center}\rule{0.5\linewidth}{0.5pt}\end{center}

\hypertarget{open-problems}{%
\subsection{9. Open Problems}\label{open-problems}}

\begin{enumerate}
\def\labelenumi{\arabic{enumi}.}
\item
  \textbf{Derive fermion masses.} The three generations are forced, but
  the mass hierarchy within and between generations is not yet derived
  from \(\ell^1\) topology.
\item
  \textbf{Derive the CKM and PMNS matrices.} The flavor mixing angles
  and CP violation phases enter as free parameters.
\item
  \textbf{Derive \(M_{\mathrm{SUSY}}\) from \(\ell^1\) criteria.}
  Eliminating \(M_{\mathrm{SUSY}}\) as a free parameter would strengthen
  the framework.
\item
  \textbf{Derive gravity from tiling defects.} The \(\{3,6\}\) tiling
  has zero Regge deficit. Curvature (gravity) requires defects. The
  mechanism by which \(\ell^1\) defects source geometric curvature is
  formalized in Paper 010.
\end{enumerate}

\begin{center}\rule{0.5\linewidth}{0.5pt}\end{center}

\hypertarget{conclusion}{%
\subsection{10. Conclusion}\label{conclusion}}

Within the \(\ell^1\) topological framework developed here, the Standard
Model gauge group emerges as the unique symmetry compatible with minimal
defect resolution, cyclic transport, and irreducibility constraints. The
generation count \(N_g = 3\) equals the number of simplex vertices. The
Weinberg angle at \(M_{\mathrm{GUT}}\) is exactly \(3/8\). The
electroweak breaking pattern is a geometric dimensional projection.

Traditional GUT groups (\(SU(5)\), \(SO(10)\), \(E_6\)) are excluded by
the transport amplitude monotonicity theorem: \(a(N)\) is strictly
decreasing in \(N\), so the triangle (\(N = 3\)) is the unique
maximizer. Within this framework, no larger group is generated or
required. It is the \(\ell^1\) minimum.

The complete derivation chain, from the monotone correction axiom
through the Standard Model gauge structure, is available in the
computational companion \texttt{gut\_l1\_derivation.py} (12 execute
functions, zero assertion failures).

\begin{center}\rule{0.5\linewidth}{0.5pt}\end{center}

\hypertarget{references}{%
\subsection{References}\label{references}}

{[}1{]} H. Georgi and S. L. Glashow, ``Unity of All Elementary-Particle
Forces,'' \emph{Phys. Rev.~Lett.} \textbf{32}, 438 (1974).

{[}2{]} H. Fritzsch and P. Minkowski, ``Unified interactions of leptons
and hadrons,'' \emph{Annals of Physics} \textbf{93}, 193 (1975).

{[}3{]} K. Abe et al.~(Super-Kamiokande Collaboration), ``Search for
proton decay via \(p \to e^+\pi^0\),'' \emph{Phys. Rev.~D} \textbf{95},
012004 (2017).

{[}4{]} Particle Data Group (R. L. Workman et al.), ``Review of Particle
Physics,'' \emph{PTEP} \textbf{2022}, 083C01 (2022). Updated 2024.

{[}5{]} J. H. Carroll, ``Paper 000: Projection Obstruction Theory:
Retraction Nonlinearity, l1 Rigidity, and Density Scaling'' (2026).

{[}6{]} J. H. Carroll, ``Paper 001: The Free \(\ell^1\) Seminorm on
Banach Presheaf Coboundaries,'' (2026).

{[}7{]} J. H. Carroll, ``Paper 002: Coordinate-Wise Additivity and the
\(\ell^1\) Norm on Finite Graph Cochains'' (2026).

{[}8{]} J. H. Carroll, ``Paper 003: Hodge Structure and Gauge
Equivalence of \(\ell^1\) Defect Fields'' (2026).

{[}11{]} J. H. Carroll, ``Paper 004: Universal Obstruction Theory - The
\(\ell^1\) Topological Framework'' (2026).

{[}12{]} J. H. Carroll, ``Paper 005: Autopoietic Cohomology: Iterative
Obstruction Repair on Causal Posets,'' (2026).

{[}13{]} J. H. Carroll, ``Paper 006: The Unification Framework:
\(\hbar\), \(c\), and \(G\) as Structural Scaling Parameters for
\(\ell^1\) Defect Systems'' (2026).

{[}14{]} J. H. Carroll, ``Paper 007: Topological Constraint on the
Electromagnetic Coupling Scale from a Discrete \(\ell^1\) Obstruction''
(2026).

\begin{center}\rule{0.5\linewidth}{0.5pt}\end{center}

\emph{Computational companion: \texttt{gut\_l1\_derivation.py} (12
execute functions, execute\_35--46, zero assertion failures)}

\begin{center}\rule{0.5\linewidth}{0.5pt}\end{center}

\hypertarget{appendix-a-classification-of-admissible-symmetry-group-on-mathbbc3}{%
\subsection{\texorpdfstring{Appendix A --- Classification of Admissible
Symmetry Group on
\(\mathbb{C}^3\)}{Appendix A --- Classification of Admissible Symmetry Group on \textbackslash mathbb\{C\}\^{}3}}\label{appendix-a-classification-of-admissible-symmetry-group-on-mathbbc3}}

\hypertarget{a.1-setup}{%
\subsubsection{A.1 Setup}\label{a.1-setup}}

Let \(V = \mathbb{C}^3\) be the complex 1-cochain space of a 2-simplex.
Assume: 1. A cyclic operator \(S \in \mathrm{End}(V)\) with \(S^3 = I\)
2. A group \(G \subset GL(3,\mathbb{C})\) acting on \(V\) 3. The
following conditions:

\hypertarget{axioms}{%
\paragraph{Axioms}\label{axioms}}

\begin{itemize}
\tightlist
\item
  \textbf{(A1) Commutation with transport}
  \[gS = Sg \quad \forall g \in G\]
\item
  \textbf{(A2) Norm preservation} \[G \subset U(3)\]
\item
  \textbf{(A3) Irreducibility} \(V\) has no nontrivial \(G\)-invariant
  subspaces
\item
  \textbf{(A4) Non-abelianity} \(G\) is non-abelian
\item
  \textbf{(A5) Connectedness} \(G\) is a connected Lie group
\end{itemize}

\hypertarget{a.2-diagonalization-of-the-shift-operator}{%
\subsubsection{A.2 Diagonalization of the Shift
Operator}\label{a.2-diagonalization-of-the-shift-operator}}

Since \(S^3 = I\), its spectrum is:
\[\{1, \omega, \omega^2\}, \quad \omega = e^{2\pi i/3}\]

Thus: \[V = V_0 \oplus V_1 \oplus V_2\] with each \(V_k\)
one-dimensional.

\hypertarget{a.3-centralizer-structure}{%
\subsubsection{A.3 Centralizer
Structure}\label{a.3-centralizer-structure}}

From (A1): \[G \subset \{T \in U(3) \mid TS = ST\}\]

Thus \(G\) lies in the \textbf{centralizer} of \(S\). Since eigenvalues
are distinct, the centralizer is:
\[Z(S) = \{ \text{diagonal matrices in eigenbasis} \}\]

\hypertarget{a.4-resolution-of-the-apparent-contradiction}{%
\subsubsection{A.4 Resolution of the Apparent
Contradiction}\label{a.4-resolution-of-the-apparent-contradiction}}

If taken alone, this would imply \(Z(S) \cong U(1)^3\), which is
abelian. But (A3) and (A4) exclude this. The key observation is that the
symmetry group is not the centralizer of a fixed \(S\), but the group
preserving the \emph{structure} of cyclic transport.

Thus, \(S\) is not fixed in a single basis; admissible transformations
act by conjugation on the cyclic structure. This lifts the restriction
\(G \not\subset Z(S)\), but instead:
\[G \subset U(3), \quad G \text{ preserves cyclic structure up to equivalence}\]

\hypertarget{a.5-classification}{%
\subsubsection{A.5 Classification}\label{a.5-classification}}

We now classify connected compact Lie subgroups of \(U(3)\) acting
irreducibly on \(\mathbb{C}^3\). By standard representation theory: *
irreducible unitary representations of dimension 3 * connected compact
groups

Possible candidates: * \(U(3)\) * \(SU(3)\) * subgroups thereof

\hypertarget{a.6-elimination}{%
\subsubsection{A.6 Elimination}\label{a.6-elimination}}

\begin{itemize}
\tightlist
\item
  \(U(3)\): contains global phase \(\implies\) violates physical
  indistinguishability
\item
  proper subgroups: either reducible, abelian, or lower dimension
\end{itemize}

\hypertarget{a.7-conclusion}{%
\subsubsection{A.7 Conclusion}\label{a.7-conclusion}}

The only connected compact Lie group acting irreducibly on
\(\mathbb{C}^3\), compatible with cyclic transport symmetry,
non-abelian, modulo global phase is: \[G = SU(3)\]

The result is not that \(SU(3)\) is assumed, but that any connected
non-abelian symmetry compatible with irreducible cyclic transport on a
three-component defect space must be \(SU(3)\) up to its center.

\end{document}
